% =====================================================================================
%  Sec. 5 -- WHAT THE CERTIFICATE CAN AND CANNOT ASSERT:
%            THE NON-VACUITY THRESHOLD OF THE LEAKAGE BOUND
%
%  NEW section; no predecessor in the v2 tree.  It resolves the live internal contradiction
%  between discussion.tex:55-58 ("demands a large fraction of the sector") and
%  results.tex:253-256 ("falls to 8 per cent"), today two opposite arguments in one paper.
%
%  This file contains TWO blocks, separated by the marked rule near the end, on the pattern
%  sec_4.tex already uses:
%    [BODY]     -> main body, after Sec. 4.  Measured: 4.1 pages at the main.tex preamble
%                  (11pt a4, margin 2.2cm, \textwidth 472.32pt), three tables included.
%                  For calibration, sec_7.tex -- also labelled "approx. 3 pp" in the plan --
%                  measures 4.0 pages in the same preamble.
%    [APPENDIX] -> after \appendix: the pipeline audit, ~1.1 pages.  NOT body text.  It
%                  carries (H1)-(H3) of the leakage bound checked against this pipeline, the
%                  ranking/tie/support declaration, and the anatomy of the displacement.
%                  If Sec. 3 ends up stating the three hypotheses with the theorem -- see the
%                  %TODO already standing in sec_3_1.tex -- this appendix becomes the pipeline
%                  CHECK of hypotheses stated there, and only its first sentence needs
%                  rewording.  Do not delete it: (H1) and (H3) are the two a referee attacks
%                  first, and neither is checked anywhere else in the manuscript.
%
%  HEADING.  The mandated heading names "Theorem B".  That name does not exist in the v3
%  numbering (the theorem is \ref{thm:leakage}, numbered by \newtheorem in main.tex), and a
%  \ref inside \section is copied into the PDF bookmarks and warns.  The heading is kept
%  verbatim except that "Theorem B" reads "the leakage bound"; the framing paragraph names
%  Theorem~\ref{thm:leakage} in its first line.  ONE-WORD INTEGRATION DECISION: if Sec. 3
%  does christen it "Theorem B" in prose, restore the original heading.
%
%  NEW labels introduced here (none exist in the v2 tree):
%     sec:frontier                (owned here; already referenced from sec_3_1.tex,
%                                  sec_4.tex, sec_6.tex, sec_7.tex)
%     sec:frontier:hyp, sec:frontier:threshold, sec:frontier:certified,
%     sec:frontier:calib, sec:frontier:obstruction
%     eq:tau, tab:thresholds, tab:shotfrontier, tab:certified
%     app:hyp  (the appendix block at the end of this file)
%
%  EXTERNAL labels used (owned elsewhere -- reconcile in one pass at integration):
%     sec:leakage    Sec. 3, the two-budget decomposition        [Sec. 3, NOT YET WRITTEN]
%     thm:leakage    Theorem B, the leakage bound                [Sec. 3, NOT YET WRITTEN]
%     eq:twobudget   the two-sided bound displayed in Sec. 3     [Sec. 3, NOT YET WRITTEN]
%     sec:retraction Sec. 3.1                                    [sec_3_1.tex]
%     sec:nogo       Sec. 6                                      [sec_6.tex]
%     sec:prices     Sec. 7                                      [sec_7.tex]
%     tab:shots      Sec. 7, the measured shot budgets           [sec_7.tex]
%     sec:physics    Sec. 8                                      [sec_8.tex]
%     fig:akwsampled, fig:scaling                                [rebuilt; MANIFIESTO F5 / F12]
%     fig:thm1iii-violation   appendix figure N3                 [MANIFIESTO N3]
%     app:counterex  Appendix B, counterexamples + pooled sweep  [NOT YET DEFINED]
%
%  tau(w_S) is DEFINED here, Eq.~\eqref{eq:tau}.  sec_6.tex currently re-derives the same
%  expression inline inside Prop.~\ref{prop:leakinv}; at integration that inline definition
%  should point here instead.  The two agree.
%
%  FIGURE THAT DOES NOT EXIST.  The plan lists a new figure N1 (the certified region in the
%  (fraction, eta) plane).  IT HAS NOT BEEN BUILT and is NOT referenced anywhere below; the
%  section is self-contained on its three tables.  See the report: it can be built from the
%  same 204-row sweep and would carry the convergence flags of Table~\ref{tab:thresholds} as
%  hatching rather than as footnote marks.  Do NOT add a \ref to it until it exists.
%
%  PROVENANCE.  Every number below is from the C3 sweep (204 evaluations: out/C3_grid.json,
%  the 54 pt_*.json point files, ties.json, validation.json, verdict_frontier.json) and the
%  four adversarial re-runs (atk2/z_*.json, z_suppK_L8.log, z_fine_L8.log, z_rank_L6.json),
%  as adjudicated in P2_FRONTERA.md.  Where a report and the adjudication disagree, the
%  adjudication is used.
%  *** %TODO(deposit), narrowed 2026-09-18 after listing release/: the leakage evaluator IS
%  deposited (src/leakage_certificate.py, src/leakage_certificate_suite.py) and so are the
%  606-subspace scans (data/cert_akw.json, cert_stress.json, cert_teqsci.json,
%  cert_frontier_6-8.json, cert_frontier_10.json, leakage_certificate_selftest.json).
%  RESOLVED 2026-09-19, and the resolution is recorded rather than the note deleted,
%  because what it used to say is the opposite of the truth and someone reading the old
%  wording could weaken the (correct) Data Availability Statement to match it.
%  The note read: "WHAT IS STILL MISSING is the 204-row frontier sweep itself --
%  C3_grid.json, the 54 pt_*.json point files, ties.json, validation.json,
%  verdict_frontier.json and the four adversarial re-runs -- i.e. the source of every
%  number in Sections 5.2 to 5.5.  This section is still the one part of the paper a
%  reader cannot regenerate."  Every item on that list is now deposited AND versioned:
%  data/c3_frontier/ holds C3_grid.json, ties.json, validation.json,
%  verdict_frontier.json, all 54 pt_*.json and adversarial/ (5 files), for 82 files
%  under git.  Sections 5.2 to 5.5 regenerate from the deposit.  Verified by listing
%  the tree and by `git ls-files data/c3_frontier/`, not by reading this comment. ***
%  *** %TODO(calc): there is NO finite-shot run at the literal configuration of this section.
%  Every subspace here is the T -> infinity limit of the declared protocol; that is stated in
%  Sec. 5.1 and must not be dropped anywhere these tables are quoted. ***
% =====================================================================================

% ======================================= [BODY] ======================================

\section{What the certificate can and cannot assert: the non-vacuity threshold of the leakage bound}
\label{sec:frontier}

Throughout this section $\mathrm{FR}^{\ast}$ is the \emph{non-vacuity threshold of the certificate of
Theorem~\ref{thm:leakage}}: the sector fraction above which the leakage branch of the bound first
improves on the trivial bound $\lVert\phi\rVert^{2}(1+w_S)$. It is a property of the inequality chain
that proves the theorem, not a frontier of the method: at $L=8$, $\eta=0.18\,t$ the acceptance
criterion of this paper, a relative $L_1$ error of $0.05$, is reached at $\mathrm{FR}\approx0.52$, while
the certificate does not reach $0.05$ until $\mathrm{FR}\approx0.98$.

Non-vacuity is a condition on one number. With $\widehat\Lambda_S=\Lambda_S/\lVert\phi\rVert$, the
leakage branch of Eq.~\eqref{eq:twobudget} improves on the trivial branch exactly when
\begin{widetext}
\begin{equation}
\frac{\widehat\Lambda_S(\eta)}{\eta}\;<\;\tau(w_S)\;:=\;\frac{1+w_S}{1+\sqrt{w_S}}-\sqrt{1-w_S},
\qquad \tau(1)=1 ,
\label{eq:tau}
\end{equation}
\end{widetext}
so a subspace holding the entire Born weight of the probe certifies something as soon as its leakage,
in units of the resolution it was asked for, drops below unity. This section is that one comparison,
made $204$ times.

\subsection{What is evaluated}
\label{sec:frontier:hyp}

We evaluate Theorem~\ref{thm:leakage} on the subspaces of Fig.~\ref{fig:scaling} at $L=6$, $8$, $10$ and
$12$, with the ranking protocol of Sec.~\ref{sec:prices}, over fractions
$0.20$--$0.95$ and $\eta=0.10$--$0.25\,t$: $204$ evaluations, of which $187$ are converged in Krylov
depth, and no row flagged not converged is quoted. To that grid we add twenty evaluations at the
published subspaces themselves, at the integer counts $246$, $2180$, $19\,183$, $124\,407$ and
$824\,504$ at $L=6$--$14$, all passing the depth test. The subspaces are the $T\to\infty$ limit of the
protocol, not finite-shot draws. What is certified is the object the code forms, the reorthogonalized
Krylov reconstruction $A_{\mathcal K}$ at a declared depth, against a reference whose drift is declared,
and no measured error within a factor $20$ of that drift is quoted. The protocol, the controls and the
costs are in Secs.~\ref{sm:frontier} and~\ref{app:hyp}.

\subsection{The threshold, measured: at every published fraction the certificate is vacuous}
\label{sec:frontier:threshold}

Evaluated on the published subspaces themselves, the
leakage branch of the bound exceeds the trivial bound $\lVert\phi\rVert^{2}(1+w_S)$ at all five sizes
and all four resolutions, by factors running from $1.66$ ($L=6$, $\eta=0.25\,t$) to $8.61$ ($L=14$,
$\eta=0.10\,t$), and by $2.40$ to $4.61$ at the $\eta=0.18\,t$ of the momentum-resolved channel;
Table~\ref{tab:published} carries the twenty values. They are monotone in both
arguments---rising with $L$ at every resolution, falling with $\eta$ at every size---so the trend runs
against us: the certificate is furthest from saying anything at the largest sector and the finest
resolution, which is where the method is aimed.

The largest classically converged case carries the evidence. At $L=12$ and the published
subspace---$\lvert S\rvert=124\,407$ of $731\,808$ determinants---$w_S=0.998622$, so
Eq.~\eqref{eq:tau} asks for $\widehat\Lambda_S/\eta<\tau(w_S)=0.9625$; the measured value is $4.236$
at $\eta=0.18\,t$, converged in depth (Sec.~\ref{sm:frontier}). Over the $606$ live subspaces of the
pooled sweep the leakage branch is the smaller of the two on only $61$. \textbf{The reconstructions at
the operating fractions of the resource scan---Fig.~\ref{fig:scaling} and the operating points of
Secs.~\ref{sec:prices} and~\ref{sec:physics}---are therefore, at present, uncertified.}
Fig.~\ref{fig:akwsampled}, drawn at $85\%$ of its sector, is the exception: there the certificate of
the depth-$500$ Krylov reconstruction is non-vacuous in all sixteen channels, at a relative bound of
$0.77$--$1.75$ against the trivial $2$, and loose, $350$--$900$ times the true error
(\texttt{cert\_akw.json}). The $L=14$ row rests on the stored ground-state checkpoint and was not
re-ranked; its caveats are in Sec.~\ref{sm:frontier}.

\begin{table*}[tp]\centering\small
\caption{\textbf{The certificate on the subspaces this paper publishes, including the size the sweep
never reached.} Each row is Theorem~\ref{thm:leakage} evaluated at the integer determinant count
$\lvert\mathcal S\rvert=\mathrm{round}(\mathrm{frac}\times\dim)$ that Fig.~\ref{fig:scaling} reports,
not at the nearest point of the grid of Table~\ref{tab:thresholds}. Columns~$5$--$8$ are the vacuity
factor---the leakage branch of Eq.~\eqref{eq:twobudget} over the trivial branch
$\lVert\phi\rVert^{2}(1+w_S)$---so a value above $1$ means the certificate says nothing the trivial
bound does not. All twenty evaluations are the $T\to\infty$ limit of the ranking protocol and pass
the depth test; controls and costs are in Sec.~\ref{app:hyp}. The last column is the measured
relative $L_1$ error at the resolution and against the threshold that \emph{define} the published
fraction ($\eta=0.15\,t$, rel-$L_1<0.05$), from code sharing neither reference, normalization nor
recursion depth with the scan that produced the fractions. At $L=14$ the count is one determinant above the deposited $824\,503$ and the ranking was not re-run; see the text. The $L=14$ row is taken at the deepest recursion run, depth $250$, where the depth indicator is $0.045$ and the factor at $\eta=0.10\,t$ still falls by $0.17\%$ between depths $200$ and $250$ ($8.627\to8.613$); the vacuity is not at risk.}\label{tab:published}
\setlength{\tabcolsep}{6pt}
\begin{tabular}{lcccccccc}
\toprule
 & & \multicolumn{2}{c}{published subspace} & \multicolumn{4}{c}{vacuity factor at $\eta/t=$} & rel-$L_1$\\
\cmidrule(lr){3-4}\cmidrule(lr){5-8}
$L$ & $\dim(N{+}1)$ & fraction & $\lvert\mathcal S\rvert$ & $0.10$ & $0.15$ & $0.18$ & $0.25$ & at $\eta=0.15\,t$\\
\midrule
$6$  & $300$          & $0.8200$ & $246$      & $4.54$ & $2.94$ & $2.40$ & $1.66$ & $0.0477$\\
$8$  & $3\,920$       & $0.5561$ & $2\,180$   & $5.65$ & $3.69$ & $3.05$ & $2.17$ & $0.0478$\\
$10$ & $52\,920$      & $0.3625$ & $19\,183$  & $6.50$ & $4.28$ & $3.54$ & $2.52$ & $0.0484$\\
$12$ & $731\,808$     & $0.1700$ & $124\,407$ & $8.04$ & $5.20$ & $4.27$ & $3.00$ & $0.0485$\\
$14$ & $10\,306\,296$ & $0.0800$ & $824\,504$ & $8.61$ & $5.59$ & $4.61$ & $3.26$ & $0.0463$\\
\bottomrule
\end{tabular}
\end{table*}

Two controls come with the table. Applying this paper's own acceptance criterion to the same sweep, with
code sharing neither reference, normalization nor recursion depth with the scan, returns the reported
fractions to under $2\%$ (Table~\ref{tab:thresholds}), and the last column of
Table~\ref{tab:published} evaluates the defining criterion at the published subspaces themselves,
$0.0463$--$0.0485$ against the threshold $0.05$. What fails is the certification of those fractions,
not the reconstruction and not the curve. In shots the gap to the non-vacuous band is a budget: reaching it multiplies the shots per
snapshot needed for the marginal determinant by $7.5$--$10$ at $L=8$ and by $20$--$35$ at $L=12$
(Table~\ref{tab:shotfrontier}).

\subsection{The strongest certified statements, and what this paper's own criterion would cost}
\label{sec:frontier:certified}

What the sweep does certify, at each size the strongest statement whose bracket is converged, is
collected in Table~\ref{tab:certified}: at $L=12$, $60\%$ of a $731\,808$-dimensional sector gives
relative $L_1\le0.717$, a factor $2.8$ below the trivial cap, tightening to $\le0.218$ at $70\%$ and
$\eta=0.25\,t$. At the acceptance criterion this paper uses everywhere
else, relative $L_1<0.05$, the certificate requires at least $0.98$ of the sector at $L=8$, the crossing
lying in $[0.977,0.983]$, and on our grid it is not reached at all at $L=6$, $10$ or $12$: at the
accuracy this paper accepts elsewhere, Theorem~\ref{thm:leakage} does not certify the operating points
of Fig.~\ref{fig:scaling}, and the shot cost of the regime where it would exceeds the saving that figure
reports.

\subsection{Where the certificate is informative, and where it is not}
\label{sec:frontier:calib}

One positive result survives the sweep. Across all sixteen $(L,\eta)$ cells the crossing of the trivial
bound occurs at a true relative $L_1$ error of $1.2\times10^{-3}$ to $7.3\times10^{-3}$, median
$2.8\times10^{-3}$: a spread of $5.9$ over true errors spanning five decades, and in every cell at least
$47$ times the drift of our own reference. The spread is a drift, not scatter, and it runs the
conservative way in $L$: the larger the system, the smaller the true error at which the certificate
first says anything. $\widehat\Lambda_S(\eta)/\eta$ is therefore a calibrated a posteriori risk
tracker, computed from $H$, $S$, $\phi$ and $\eta$ alone, with four limits (Sec.~\ref{sm:frontier}).
It locates one error level per cell and is not an error estimate. The comparison needs $\tau(w_S)$,
which a user can discharge only conservatively: an assumed floor $w_S\ge0.995$ gives $\tau\ge0.9280$, at
the cost of at most $7\%$ in the threshold. Sixteen cells are the whole of the evidence, and no
extrapolation to $L=14$ is made. And a two-constant rule in $1-w_S$, fitted on some sizes and tested
on the others, localizes the crossing error more tightly than the certificate on eleven of the
fourteen train/test splits (\texttt{null\_ws7.py}), so we claim no superiority for the certificate as a
tracker against the captured weight. What distinguishes it is that it is an inequality with a proof,
whose decision survives replacing $w_S$ by a floor, while the fitted rule asks for $1-w_S$ to
$10^{-5}$--$10^{-4}$ and, stripped of its second parameter, is the looser of the two on thirteen of the
fourteen splits.

\subsection{Where the threshold comes from}
\label{sec:frontier:obstruction}

Measured rung by rung along the chain of inequalities that proves Theorem~\ref{thm:leakage}, at $L=8$,
the displacement of the threshold is $44\%$ the pointwise step to a Hilbert-space distance, $38\%$ the
operator-norm step $\lVert R(z)\rVert\le1/\eta$, $11\%$ the Feshbach/Minkowski split and $4.6\%$
Cauchy--Schwarz in $\omega$, whose total margin for tightening is a factor $1.05$ typically and $3.29$
at worst. The subspace adds one structural fact. For a subspace holding the whole probe, $\Lambda_S=0$
requires $\operatorname{span}(S)\supseteq\mathcal K(H,\phi)$, and the reflection $j\to-j$ about the
seed site commutes with $H$ and has $\phi$ as an eigenvector, so every vector of $\mathcal K$ vanishes on
the reflection-fixed determinants of the opposite parity---$4$, $18$, $48$, $200$ and $600$ at
$L=6$--$14$---while a reflection-projected power iteration (\texttt{z\_suppK\_sym.py}) finds every other
determinant in the support: $296$, $3902$, $52\,872$, $731\,608$ and $10\,305\,696$ at $L=6$--$14$. At
full captured weight, therefore, no coordinate subspace short of all of them has zero leakage. That is a
statement that $\Lambda_S\neq0$, not about its size; the vacuity at the published fractions, all below
full captured weight, is the measurement of Table~\ref{tab:published} and not a deduction from it. Nor
is the ranking the weak link: a Krylov-diagonal oracle ranking gives a \emph{worse} certificate than the
Born ranking at $L=6$, $\eta=0.18\,t$, at every fraction checked from $0.50$ on. At full captured weight
the leakage branch is of first order in the boundary coupling and the error of second, and the product
$\Lambda_P\tilde\Lambda_P/\eta^{2}$ is itself an a posteriori certificate of second order
(Sec.~\ref{sm:frontier}); every published fraction lies below full captured weight, and whether a
certificate of second order exists there is open. Finally, for the momentum-resolved probe of
Fig.~\ref{fig:akwsampled} the Krylov space lies inside one momentum block of exactly $\dim(N{+}1)/L$
determinants, on which the certificate becomes non-vacuous at $48$--$50$ of $50$ determinants ($L=6$)
and $455$--$483$ of $490$ ($L=8$) (\texttt{z\_kblock.py}), a factor $5$--$7$ fewer than the determinant
basis requires; the gain is $1/L$ against an exponential, and whether the support fills the block at
$L\ge10$ is unmeasured.

